\scanpage{169}
\mspage{24}
pour ne devenir plus qu'un pâle élément d'une ``catégorie homotopique'' (ou \unsure{pro-homotopique|prohomotopique}), déduite de celle des topos par un processus de localisation à \unsure{brin|bras|brins} de zinc -- à première vue, tout ce qui reste encore à ce pauvre $X$, ce sont ses $\pi_i$ -- et ses groupes de \unsure{cohomologies|cohomologie} à \unsure{coefficient|coefficients} constants -- ou à la rigueur constant tordus. Quand on fouille un peu plus dans cette direction ``que reste-t-il encore à ce pauvre $X$'' on tombe précisément sur les \inserted{$n$-}champs \uline{loc.} \uline{constants} (car un $X\mathrel{\unsure{\tikz[baseline=-.5ex]{\draw (0,0)--(.52,0); \node[anchor=base,inner sep=0pt] at (.42,-.08) {$f$};}|\overset{f}{\longrightarrow}}}X'$ qui est une équivalence d'homotopie induit une $(n+1)$-équivalence entre les cat. de $n$-champs loc. \unsure{constant|constants} sur $X$ et sur $X'$) -- qui contiennent bien sûr les complexes \inserted{de \unsure{chaines|chaînes}} abéliens \inserted{de long. $n$} à faisceaux de coh. loc. constants, et l'hypercohomologie de ceux-ci. C'est ainsi qu'on arrive à \deleted{\unsure{la notion|la notion du}} \inserted{ce} triangle de \inserted{\uline{\unsure{\textquoteleft|``}Objets\unsure{\textquoteright|''|{}}}} se déterminant mutuellement
\diagram{\begin{tikzcd}[ampersand replacement=\&,column sep=-1.25em,row sep=large]
\& \begin{gathered}
\text{topos (ou espace top., ou}\\
\text{ens. }{}^{1}\!/\!{}_{2}\,\text{\unsure{simpl|simpl.}) }\text{\uuline{\unsure{modulo|modúlo}}}\\
n\text{-}\text{\uline{homotopie}}
\end{gathered} \arrow[dl,leftrightarrow] \arrow[dr,leftrightarrow] \& \\
\begin{gathered}
n\text{-groupoïdes}\\
\text{(à }n\text{-équiv. \unsure{pr.|p.|près})}
\end{gathered} \arrow[rr,leftrightarrow] \& \& \begin{gathered}
n\text{-catégories ``spéciales'' (à}\\
n\text{\unsure{..|-}équivalence près)}
\end{gathered}
\end{tikzcd}}
% archive sheet number: 169
